% Owner B: Siyuan Jiang. Sequence-level design is not yet finalized.
\subsection{Construct Strategy and Materials}
The starting outline proposes an N-terminal EGFP--TFEB fusion expressed from a CMV promoter, using pEGFP-C1 as a candidate backbone. The final design will specify a single backbone and a versioned TFEB coding sequence, then verify the promoter, fusion junction, multiple cloning site, bacterial replication origin, and bacterial and mammalian selection markers. The original outline's fusion orientation remains provisional until the sequence and literature checks are complete.

For an EGFP--TFEB construct with TFEB at the C terminus, the TFEB termination codon should be retained unless an additional downstream fusion is explicitly intended. For the alternative TFEB--EGFP orientation, the native TFEB termination codon must be omitted to permit translation into EGFP. Removing a stop codon does not by itself establish the correct reading frame; both junctions require sequence-level verification.

The planned molecular materials are a defined human TFEB template, the selected EGFP expression vector, high-fidelity DNA polymerase, two compatible restriction enzymes, DNA ligase, competent \textit{Escherichia coli} DH5$\alpha$, and media containing the antibiotic specified by the verified backbone. A circular plasmid map will show the final sequence features and cloning junctions.

\subsection{Sequence Retrieval and Primer Design}
A human TFEB reference transcript and coding sequence will be retrieved from NCBI with accession and version recorded. Each cloning primer will be reported in the 5$'$ to 3$'$ direction and annotated by protective bases, restriction site, any frame-adjustment or linker bases, and template-annealing sequence. The selected enzyme pair must cut the intended backbone at the required positions without cutting the TFEB insert internally. Primer length, melting temperature, and compatibility with the amplification conditions will be documented after the sequence is fixed.

\subsection{PCR Amplification and Restriction Cloning}
The TFEB insert will be amplified with a high-fidelity polymerase. A no-template control and an appropriate amplification control will accompany the reaction. The intended amplicon will be purified, and both insert and vector will be digested and purified before ligation. An initial insert-to-vector molar ratio of approximately 3:1 is proposed. The ligation will be transformed into competent DH5$\alpha$ cells, followed by selection using the backbone's confirmed antibiotic marker. The reaction composition, thermal program, digestion conditions, and transformation details remain to be specified for the selected reagents.

\subsection{Clone Screening and Sequence Verification}
Candidate colonies will be screened by colony PCR and diagnostic restriction digestion. Plasmid DNA from positive clones will be purified for Sanger sequencing with vector-flanking and, where necessary, internal primers. A construct will be accepted only when the complete TFEB insert, both junctions, fusion reading frame, and terminal stop codon match the intended sequence without unintended substitutions, insertions, or deletions. An expected-size band alone will not establish sequence identity.
